\documentclass[12pt,a4paper]{article}
\usepackage[utf8]{inputenc}
\usepackage[T1]{fontenc}
\usepackage[brazil]{babel}
\usepackage{amsmath,amssymb}
\usepackage{geometry}
\geometry{margin=2.5cm}
\title{Material de Estudo: Din\^amica Molecular I}
\author{Princ\'ipio de Modelagem Matem\'atica}
\date{Unidade 27}
\begin{document}
\maketitle

\section*{Objetivo}
Compreender os fundamentos da Din\^amica Molecular (DM): potencial de Lennard-Jones, algoritmos de Verlet, e condi\c{c}\~oes peri\'odicas de contorno.

\section*{Fundamentos da DM}

A DM simula o movimento cl\'assico de $N$ \'atomos resolvendo as equa\c{c}\~oes de Newton:
\[
m_i\ddot{\vec{r}}_i = \vec{F}_i = -\nabla_i V(\{\vec{r}\}),
\]
onde $V$ \'e a energia potencial total do sistema.

\section*{Potencial de Lennard-Jones}

Descreve a intera\c{c}\~ao entre dois \'atomos n\~ao ligados:
\[
V_{LJ}(r) = 4\varepsilon\!\left[\left(\frac{\sigma}{r}\right)^{12} - \left(\frac{\sigma}{r}\right)^6\right].
\]
\begin{itemize}
  \item Termo $r^{-12}$: repuls\~ao de Pauli (duas part\'iculas n\~ao se sobrepõem).
  \item Termo $r^{-6}$: atra\c{c}\~ao de van der Waals (dipolo induzido).
  \item M\'inimo em $r_\text{eq} = 2^{1/6}\sigma$ com $V = -\varepsilon$.
  \item For\c{c}a: $F(r) = -dV/dr = 4\varepsilon[12\sigma^{12}/r^{13} - 6\sigma^6/r^7]$.
\end{itemize}

\section*{Algoritmo de Verlet}

\[
\vec{r}(t+\Delta t) = 2\vec{r}(t) - \vec{r}(t-\Delta t) + \Delta t^2 \vec{a}(t).
\]
\begin{itemize}
  \item Erro por passo: $O(\Delta t^4)$ para posi\c{c}\~oes.
  \item Conserva energia muito bem em longo prazo (simpl\'etico).
  \item Problema: velocidade n\~ao est\'a expl\'icita.
\end{itemize}

\subsection*{Velocity-Verlet}
\[
\vec{r}(t+\Delta t) = \vec{r}(t) + \Delta t\vec{v}(t) + \frac{\Delta t^2}{2}\vec{a}(t).
\]
\[
\vec{v}(t+\Delta t) = \vec{v}(t) + \frac{\Delta t}{2}[\vec{a}(t) + \vec{a}(t+\Delta t)].
\]
Posi\c{c}\~oes e velocidades s\~ao atualizadas explicitamente.

\section*{Condi\c{c}\~oes Peri\'odicas de Contorno (PBC)}

Para eliminar efeitos de superf\'icie e simular um sistema ``infinito'':
\begin{itemize}
  \item A caixa de simula\c{c}\~ao \'e replicada em todas as dire\c{c}\~oes.
  \item Uma part\'icula que sai pela direita entra pela esquerda.
  \item \textbf{Conven\c{c}\~ao de m\'inima imagem:} usa-se a c\'opia mais pr\'oxima de cada par.
\end{itemize}

\section*{Exemplos Resolvidos}

\subsection*{Exemplo 1 --- Equil\'ibrio do LJ}
Para $\varepsilon=1$ e $\sigma=1$ (unidades reduzidas): $r_\text{eq} = 2^{1/6} \approx 1{,}12$. Energia m\'inima: $V(r_\text{eq}) = -1$ (em unidades de $\varepsilon$).

\subsection*{Exemplo 2 --- Verlet: 1 passo}
Part\'icula 1D: $x(0)=0$, $v(0)=1$ \AA/ps, $a(0)=-2$ \AA/ps$^2$, $\Delta t=0{,}5$ ps.
$x(0{,}5) = 0 + 0{,}5\times1 + 0{,}25\times(-2)/2 = 0{,}5 - 0{,}25 = 0{,}25$ \AA.
$v(0{,}5) = 1 + 0{,}5/2\times(-2 + a(0{,}5))$ (precisa de $a(0{,}5)$).

\section*{Exerc\'icios}

\begin{enumerate}
  \item \textbf{(LJ)} Em unidades reduzidas ($\varepsilon=\sigma=1$): (a) calcule $V(r)$ para $r=0{,}9, 1{,}0, 1{,}12, 1{,}5, 2{,}0$. (b) Trace esbo\c{c}o da curva. (c) Em que $r$ a for\c{c}a \'e zero?
  
  \item \textbf{(For\c{c}a LJ)} Calcule a for\c{c}a $F(r) = -dV/dr$ para $r=1{,}0$ e $r=1{,}12$ (em unidades reduzidas). Interprete o sinal.
  
  \item \textbf{(Velocity-Verlet)} Uma part\'icula em 1D: $x_0=0$, $v_0=2$ m/s, for\c{c}a constante $F=-1$ N/kg. Aplique Velocity-Verlet com $\Delta t=0{,}1$ s por 5 passos. Compare com a solu\c{c}\~ao anal\'itica.
  
  \item \textbf{(PBC)} Uma caixa 1D de tamanho $L=10$ tem uma part\'icula em $x=9{,}5$ que d\'a um passo de $+2$. Para onde vai? Escreva a regra geral: $x_\text{novo} = x + \Delta x - L\lfloor(x+\Delta x)/L\rfloor$.
  
  \item \textbf{(Cutoff)} Na pr\'atica, o potencial LJ \'e truncado em $r_c = 2{,}5\sigma$. (a) Qual a energia do LJ em $r_c$? (b) Por que truncar reduz o custo computacional de $O(N^2)$?
  
  \item \textbf{(Unidades reduzidas)} Na DM com LJ, \'e comum usar unidades reduzidas: $r^* = r/\sigma$, $T^* = k_BT/\varepsilon$, $t^* = t/(\sigma\sqrt{m/\varepsilon})$. Converta: argônio tem $\sigma=3{,}4$\,\AA, $\varepsilon/k_B=120$\,K, $m=6{,}63\times10^{-26}$\,kg. Qual \'e a unidade de tempo em ps?
  
  \item \textbf{(Conserva\c{c}\~ao)} Num sistema isolado (NVE), a energia total $E = K + V$ deve ser conservada. Se numa simula\c{c}\~ao $K$ e $V$ oscilam mas $E$ \'e constante, o que isso indica? Se $E$ cresce lentamente, qual \'e o prov\'avel problema?
\end{enumerate}

\section*{Resumo R\'apido}
\begin{itemize}
  \item LJ: repuls\~ao $r^{-12}$ + atra\c{c}\~ao $r^{-6}$; m\'inimo em $r=2^{1/6}\sigma$.
  \item Velocity-Verlet: posi\c{c}\~ao e velocidade expl\'icitas; erro $O(\Delta t^4)$.
  \item PBC: elimina efeitos de borda; conven\c{c}\~ao de m\'inima imagem.
\end{itemize}

\section*{Refer\^encias}
\begin{itemize}
  \item Frenkel, D.; Smit, B.\ \textit{Understanding Molecular Simulation}. Academic Press, 2002.
  \item Allen, M.\ P.; Tildesley, D.\ J.\ \textit{Computer Simulation of Liquids}. Oxford, 2017.
  \item Notas de aula da disciplina.
\end{itemize}

\end{document}
